
\chapter{Fazit und Ausblick}
\label{chap:fazit_und_ausblick}

\section{RQ1 beantworten}

\section{RQ2 beantworten}

\section{RQ3 beantworten}

\section{RQ4 beantworten}

\section{RQ0 beantworten}


% \section{Fazit}
% Das Ziel dieser Arbeit war die Entwicklung sowie Evaluierung eines datengetriebenen Systems zur automatisierten Zuordnung von EPLAN-Artikeln \parencite{EPLANArtikelverwaltung} zu EPLAN-Makros \parencite{EPLAN-Makros}.
% Die zugrunde liegende Problemstellung wurde als eine Learning-to-Rank-Aufgabe formuliert, bei der das Modell lernen soll, die relevantesten Ziel-Makros \parencite{EPLAN-Makros} für einen gegebenen EPLAN-Artikel \parencite{EPLANArtikelverwaltung} zu priorisieren.
% Dabei soll eine Rangliste möglicher Ziel-Makros erzeugt werden.


% Hierzu wurde eine End-to-End-Pipeline entwickelt, welche die Datenaufbereitung, die Erzeugung von Kandidatenprofilen, die Berechnung strukturierter Anfrage-Kandidat-Merkmale, 
% das Training eines LightGBM LambdaRank-Modells sowie die automatische Evaluation umfasst.
% Die Evaluierung der entwickelten Lösung erfolgte anhand eines realen Datensatzes von EPLAN-Artikeln \parencite{EPLANArtikelverwaltung} und den zugehörigen EPLAN-Makros \parencite{EPLAN-Makros}. 


% Die Datenbasis umfasst 1.559 EPLAN-Artikel \parencite{EPLANArtikelverwaltung}, 375 Ziel-Makros \parencite{EPLAN-Makros} und 55 Makrofamilien.
% Aus diesen Daten wurden anschließend 375 Kandidatenprofile erzeugt und durch 2.381 technische, lexikalische und strukturierte Merkmale beschrieben. 
% Die Trainingsmatrix umfasst 369.000 Anfrage-Kandidat-Paare und bildet damit die Grundlage für das Training des LightGBM-LambdaRank-Modells.


% Es wurden neun unterschiedliche Modellvarianten im Rahmen einer systematischen Konfigurationsstudie untersucht. 
% Dabei wurden verschiedene Familiengewichtungen, Labeldefinitionen, Profilbreiten und Regularisierungsstufen bewertet.
% Die beste Leistung wurde dabei durch die Konfiguration A08 erzielt, welche einen Target Recall@1 von 34,78\,\%, einen Family Recall@1 von 87,35\,\%, einen Target Recall@10 von 67,98\,\% sowie einen Mean Reciprocal Rank von 45,34\,\% aufweist.


% Die Ergebnisse zeigen, dass das entwickelte System in der Lage ist, die korrekte Makrofamilie in 8 bis 9 von 10 Fällen korrekt zu erkennen.
% Die Analyse der vorhandenen Fehler zeigte gleichzeitig, dass die größere Herausforderung in der präzisen Erkennung der Variantenunterscheidung innerhalb derselben Makrofamilie liegt.
% Mehr als die Hälfte aller Fehlklassifikationen entsteht dabei innerhalb der korrekten Familie, während fachlich unverbundene Familien nur einen geringen Anteil der Fehler verursachen.


% Die Häufigkeiten der Ziel-Makros folgen dabei einer sogenannten Long-Tail-Verteilung. Darunter wird eine stark ungleichmäßige Verteilung verstanden. Dabei kommen wenige Ziel-Makros sehr häufig vor, während viele weitere Ziel-Makros nur selten oder einmalig vorkommen. 
% Die durchgeführten Analysen zeigen außerdem, dass sich die Leistungsgrenzen des entwickelten Systems primär durch diese Long-Tail-Verteilung mit 213 Singleton-Targets und 279 Ziel-Makros \parencite{EPLAN-Makros} mit weniger als drei Trainingsbeispielen ergeben.
% Diese erschweren den Aufbau robuster Kandidatenprofile. Die zusätzliche Leave-One-Out-Analyse zeigt darüber hinaus, dass die ursprünglich vermutete Generalisierungslücke durch die Konstruktion der Kandidatenprofile verursacht wurde.
% Demgegenüber verhält sich die Generalisierung zwischen Validierungs- und Testdaten stabil und liefert keine Hinweise auf ein kritisches Overfitting des verwendeten LightGBM-LambdaRank-Modells.


% Zusammenfassend konnte nachgewiesen werden, dass der entwickelte Learning-to-Rank-Ansatz eine geeignete Grundlage für die automatisierte Zuordnung von EPLAN-Artikeln \parencite{EPLANArtikelverwaltung} zu EPLAN-Makros \parencite{EPLAN-Makros} darstellt.

% Die künftigen Leistungssteigerungen sind vielmehr durch eine verbesserte Datenbasis als durch zusätzliche Hyperparameter-Optimierungen zu erwarten. 
% Zusätzliche Variantenbeispiele sowie spezialisierte Verfahren zur Variantenauflösung innerhalb der Makrofamilien können die zukünftige Leistungsfähigkeit des Systems weiter verbessern.

% Somit kann die Forschungsfrage trotz der identifizierten Limitationen und Herausforderungen des Systems als positiv beantwortet werden, da nachgewiesen werden konnte,
% dass der limitierende Faktor nicht die Modellarchitektur, sondern die Datenbasis darstellt.


% \section{Ausblick}

% Die Ergebnisse dieser Arbeit zeigen mehrere Ansatzpunkte für künftige Weiterentwicklungen auf.

% Die größte erwartbare Leistungssteigerung liegt in der Erweiterung der vorhandenen Datenbasis. Besonders für seltene Ziel-Makros \parencite{EPLAN-Makros} sollte eine größere Basis an Trainingsbeispielen geschaffen werden.
% Dadurch lässt sich die heutige Dominanz von Singleton-Profilen reduzieren und der Raum der Kandidaten-Profile würde eine stabilere Struktur aufweisen.
% Dies würde dazu führen, dass die breitere Datenbasis dem Modell mehr Informationen für eine bessere Generalisierung zur Verfügung stellt.


% Ein weiterer Optimierungspunkt liegt in der Verbesserung der Variantenauflösung innerhalb einer Makrofamilie. 
% Hierfür sind detailliertere Artikelmerkmale erforderlich, die es ermöglichen, die einzelnen Artikel präziser voneinander abzugrenzen.
% Des Weiteren wäre der Ansatz einer mehrstufigen Architektur mit einem globalen Modell zur Familienerkennung und einem weiteren separaten Modell zur Variantenauflösung innerhalb der Makrofamilie denkbar.
% Dadurch könnten dann auch eventuelle Fehlerfälle zwischen der Familienerkennung und der Variantenauflösung systematischer analysiert und kompensiert werden.


% Weiterführend stellt die Extraktion von Merkmalen auch einen interessanten Ansatz dar. Seltene oder redundant genutzte Merkmale könnten gezielt identifiziert und reduziert werden. 
% Dadurch würde die Modellkomplexität verringert und die Generalisierungsfähigkeit des Modells verbessert werden.


% Eine weitergehende Idee für die Zukunft wäre die Integration von zusätzlichem Domänenwissen von Experten oder externen Wissensquellen, welche zusätzliche technische Merkmale, aber auch semantische Zusammenhänge zwischen den Artikeln bereitstellen könnten.
% Diese Kombination aus Daten und Metadaten könnte einen vielversprechenden Ansatz für die Modellierung von spezifischeren Variantenunterscheidungen darstellen.
